%% Resumo
\begin{resumo}
Este Projeto de Conclusão de Curso apresenta a concepção e a implementação de um Sistema Web Administrativo voltado
para a confeitaria Bom Gosto, visando substituir métodos manuais de controle por uma gestão digital centralizada. 
A pesquisa adota uma metodologia aplicada e descritiva, utilizando a modelagem de sistemas com a Linguagem de Modelagem Unificada (UML)
para mapear o fluxo de negócios, além do desenvolvimento modular dividido entre a gestão administrativa de pedidos 
e a gestão de produção. Como resultados, o sistema viabiliza o registro e o acompanhamento detalhado de encomendas e
itens de pronta-entrega, centraliza o agendamento de retiradas, automatiza a baixa de matérias-primas no estoque e 
calcula com precisão o custo de fabricação de cada receita. Conclui-se que a automação proposta elimina a dependência
de anotações informais, reduz os riscos de atrasos e desperdícios de ingredientes, melhora a organização das rotinas diárias
e proporciona maior segurança e eficiência operacional à administração do estabelecimento.
\vspace{\onelineskip}
\noindent

\textbf{Palavras-chaves}: Automação comercial. Gestão administrativa. Confeitaria. Sistema web. Eficiência operacional.

\end{resumo}